% Cote 119: offprints, reports, courses, CV and letters (typed and printed part): archive cote 119, site page 206 (PDF page 207).
% typescript. Draft compared: none.
% AI-transcribed by Claude Opus 5.5 (claude-opus-5-5, adaptive thinking, effort high), Message Batches API, batch r119_full (msgbatch_01DwZNm5Y7YsKBWyAN4sFs8w), 2026-11-13; not yet checked by eye.
\scanpage{119/206}
\mspage{- 11 -}

c'est exposé dans [12] et le travail de A. Borel et J.P. Serre [*], dans un
contexte plus ambitieux cela donne l'imposant séminaire [16, SGA 7]. Dans le
même esprit, cf. [12 bis].

Par ailleurs, l'idée (que je semble avoir été le premier à introduire
avec ma formulation du théorème de Riemann-Roch) de reformuler un théorème sur
une variété (dû en l'occurrence à F. Hirzebruch) en un théorème plus général
sur un morphisme de variétés, a connu par la suite une grande fortune, non
seulement en géométrie algébrique, mais aussi en topologie algébrique et topo-%
logie différentielle (à commencer par la ``formule de Riemann-Roch différentiable'',
développée par M.F. Atiyah et F. Hirzebruch sous l'inspiration de ma formulation
``relative'' du théorème de Riemann-Roch).

h) \uline{Schémas abéliens} :

En termes plus classiques, ce sont les familles de variétés abéliennes,
paramétrées par un schéma quelconque. Les résultats les plus importants que
j'y ai établis sont le ``\uline{théorème de réduction semi-stable}'' et ses conséquen-%
ces et variantes [16, SGA 7, Exp. IX], le théorème d'\uline{existence de morphismes}
\uline{de schémas abéliens} contenu dans [17] et ses variantes (généralisé par P. Deligne
[*] en un théorème sur la cohomologie de Hodge-De Rham relative d'une famille
de variétés projectives complexes non singulières), enfin une théorie des \uline{défor-}%
\uline{mations infinitésimales des schémas abéliens} (non publiée sur une base quelcon-%
que), en termes de la déformation d'une filtration de Hodge sur un $H^1$ relatif
de De Rham (interprété comme une cohomologie cristalline).

i) \uline{Groupes de monodromie}:

Mes principales contributions sont exposées (en partie par P. Deligne)
dans le premier volume de [16, SGA 7], donnant des propriétés fondamentales de
l'action du groupe de monodromie locale sur la cohomologie comme sur le groupe
fondamental d'une fibre. Parmi les principales applications, il y a le théorème
de ``réduction semi-stable'' des schémas abéliens signalé au paragraphe précédent.

j) \uline{Divagations motiviques}:

Nous entrons ici dans le domaine du rêve éveillé mathématique, où on
s'essaie à deviner ``ce qui pourrait être'', en étant aussi insensément optimiste
que nous le permettent les connaissances parcellaires que nous avons sur les
propriétés arithmétiques de la cohomologie des variétés algébriques. La notion
de motif peut se définir en toute rigueur avec les moyens du bord (c'est fait par

\margin{206}
